\chapter*{Abstract}

This thesis investigates how the confidence of a simple machine-learning model
relates to the out-of-sample performance of a rule-based trading strategy on
the gold market. The study is motivated by the observation that rule-based trading
strategies are transparent and interpretable, but their empirical
performance is often mixed, while more complex machine-learning approaches may
increase the risk of overfitting. Against this background, the thesis examines
whether a simpler and more interpretable machine-learning method can be used to
distinguish between lower-quality and higher-quality trade signals without
introducing unnecessary complexity.

The study adopts an experimental research strategy combined with quantitative
analysis of secondary data. Historical hourly price data for the broker symbol
XAUUSD is retrieved from MetaTrader 5 and used to construct a baseline
rule-based trading strategy based on technical indicators. A decision tree
classifier is then trained to evaluate candidate trade signals generated by the
baseline strategy. For each candidate signal, the model produces a confidence
score, operationalized as the predicted probability that the signal will lead
to a favorable realized outcome under the predefined exit logic. These
confidence scores are divided into ordered ranges, and the final strategy is
evaluated through out-of-sample backtesting in a deterministic Python pipeline
using MetaTrader 5 market data.

The analysis compares quantitative performance measures across the confidence
ranges, including win rate, average trade return, profit factor, Sharpe ratio,
maximum drawdown, and number of trades. By focusing on decision-tree-based
signal evaluation rather than full predictive modelling, the thesis examines
whether a simple and interpretable machine-learning model can identify more
favorable trading conditions while preserving transparency and limiting
overfitting risk.
The final three-year H1 experiment showed that the strongest confidence band
outperformed the unfiltered multi-indicator baseline out of sample, although
the relationship between confidence and performance was not strictly monotonic
across all five ranges.
\par\medskip
\thispagestyle{empty}
%The abstract and the keywords should not exceed the limit of this page

\emph{Keywords:} algorithmic trading, rule-based trading, machine learning, decision tree, signal evaluation
